\section{查找}

\begin{frame}
    \frametitle{\insertsectionhead：引例}
    \begin{exampleblock}{查询图书编号}
        一批图书编号为$10,18,21,25,32$。需要判断编号$25$是否存在，
        有时还要列出编号在$18$至$25$之间的图书。
        如果查询频繁，应如何组织和查找这些记录？
    \end{exampleblock}
    \begin{block}{关键需求}
        精确查找与范围查询的要求不同，还要考虑数据是否有序以及更新代价。
    \end{block}
\end{frame}

\begin{frame}
    \frametitle{\insertsectionhead}
    \begin{block}{查找(Search)}
        \begin{itemize}
            \item 根据给定键定位记录，或判断记录不存在
            \item 顺序查找依次比较元素；无需有序，最坏时间为$O(n)$
            \item 查找算法的选择取决于数据是否有序、存储方式及更新频率
        \end{itemize}
    \end{block}
\end{frame}

\begin{frame}
    \frametitle{\insertsectionhead}
    \begin{block}{二分查找(Binary Search)}
        \begin{itemize}
            \item 前提：序列按键有序，且可高效随机访问
            \item 每次比较中间元素，将候选区间缩小至一半
            \item 以下使用左闭右开区间$[l,r)$，未找到时返回$-1$
        \end{itemize}
    \end{block}
\end{frame}

\begin{fragile}
    \frametitle{\insertsectionhead}
    \begin{algorithmx}[alg:binary-search]{二分查找}
        \Procedure{BinarySearch}{$s,n,key$}
        \State $l\gets 0,\ r\gets n$
        \While{$l<r$}
        \State $m\gets l+\lfloor(r-l)/2\rfloor$
        \If{$s_m=key$}
        \State \textbf{return} $m$
        \ElsIf{$s_m<key$}
        \State $l\gets m+1$
        \Else
        \State $r\gets m$
        \EndIf
        \EndWhile
        \State \textbf{return} $-1$
        \EndProcedure
    \end{algorithmx}
\end{fragile}

\begin{frame}
    \frametitle{\insertsectionhead}
    \begin{exampleblock}{二分查找实例}
        在$(10,18,21,25,32)$中查找$25$：
        \[
            [0,5)\xrightarrow{21<25}[3,5)
            \xrightarrow{32>25}[3,4)
            \xrightarrow{25=25}\text{找到下标}3
        \]
    \end{exampleblock}
    \begin{itemize}
        \item 最好时间$O(1)$，最坏时间$O(\log n)$；迭代实现辅助空间$O(1)$
        \item 键重复时，该算法返回任意匹配位置；首末位置需调整边界规则
        \item 空序列与查找失败时，区间为空，返回$-1$
    \end{itemize}
\end{frame}

\begin{frame}
    \frametitle{\insertsectionhead}
    \begin{block}{散列表(Hash Table)}
        \begin{itemize}
            \item 散列函数将键映射到表中位置；不同键可能映射到同一位置，称为冲突
            \item 冲突可用链地址法或开放定址法处理；查找和插入需使用一致的规则
            \item 在散列分布良好且装载因子受控时，查找平均为$O(1)$；最坏可为$O(n)$
            \item 普通散列表不保持键的顺序，不适合直接进行范围查询
        \end{itemize}
    \end{block}
    \begin{exampleblock}{链地址法示例}
        表长$5$，$h(k)=k\bmod 5$：键$12$与$17$均落在桶$2$，
        在该桶中依次比较即可区分。
    \end{exampleblock}
\end{frame}

\begin{frame}
    \frametitle{\insertsectionhead}
    \begin{block}{散列函数与装载因子}
        \begin{itemize}
            \item 散列函数$h(k)$把键$k$映射到表下标，计算应快速且分布尽量均匀
            \item 除留余数法示例：表长$m=7$，$h(k)=k\bmod 7$
            \item 已存储$n$个键时，装载因子$\alpha=n/m$反映表的拥挤程度
            \item 装载因子增加通常会增加冲突；可扩容并重新散列所有键
        \end{itemize}
    \end{block}
\end{frame}

\begin{frame}
    \frametitle{\insertsectionhead}
    \begin{block}{链地址法(Separate Chaining)}
        \begin{itemize}
            \item 每个桶保存映射到该位置的键；冲突键放入同一链表
            \item 插入、查找和删除先定位桶，再在该桶中查找目标键
            \item 均匀散列假设下平均查找时间为$O(1+\alpha)$；最坏为$O(n)$
            \item 桶中可存多个键，因此$\alpha$可以大于$1$
        \end{itemize}
    \end{block}
    \begin{exampleblock}{例：$m=7$，$h(k)=k\bmod 7$}
        插入$10,17,24$后，桶$3$的链表可为$10\rightarrow17\rightarrow24$。
        查找$17$需在桶$3$中比较，查找$31$遍历该桶后失败。
    \end{exampleblock}
\end{frame}

\begin{frame}
    \frametitle{\insertsectionhead}
    \begin{block}{开放定址法(Open Addressing)：线性探测}
        \begin{itemize}
            \item 全部键存于表内；冲突时依次探测$(h(k)+i)\bmod m$
            \item 插入时选第一个可用位置；查找沿相同探测序列进行
            \item 找到目标或遇到从未使用过的空槽时结束；最多探测$m$次
        \end{itemize}
    \end{block}
    \begin{exampleblock}{例：$m=7$，$h(k)=k\bmod 7$}
        依次插入$10,17,24$，分别存入下标$3,4,5$；查找$24$依次检查$3,4,5$。
    \end{exampleblock}
\end{frame}

\begin{frame}
    \frametitle{\insertsectionhead}
    \begin{block}{删除与性能}
        \begin{itemize}
            \item 开放定址删除时须留下\textbf{删除标记}，不能直接改为从未使用的空槽，
                  否则可能截断后续键的查找路径
            \item 插入可复用删除槽，但若需避免重复键，应继续探测至找到该键或确认不存在
            \item 开放定址的$\alpha<1$；表接近满时探测会变长，需及时扩容并重建
            \item 在良好的散列与受控装载因子下，平均查找可为$O(1)$；最坏为$O(n)$
        \end{itemize}
    \end{block}
    \begin{alertblock}{思考}
        上例删除$17$后，为什么不能把下标$4$直接标为从未使用？
    \end{alertblock}
    \begin{block}<2-|handout:0>{参考答案}
        查找$24$会在下标$4$提前结束，无法到达下标$5$。
    \end{block}
\end{frame}


\begin{frame}
    \frametitle{\insertsectionhead：回到图书编号查询}
    \begin{exampleblock}{选择查找方式}
        \begin{itemize}
            \item 有序数组中二分查找$25$：先比较$21$，再比较$32$，最后找到$25$
            \item 表长为$5$、$h(k)=k\bmod5$的链地址散列表中，
                  $10$和$25$同在桶$0$；查找$25$须在该桶中比较
            \item 范围$[18,25]$：在有序数组中找到起点并顺序扫描，
                  输出$18,21,25$；普通散列表不保持顺序
        \end{itemize}
    \end{exampleblock}
    有序数组适合随机访问，但插入删除可能移动元素；散列表适合精确查找。
\end{frame}

\begin{frame}{查找：思考与自测}
    \begin{itemize}
        \item 在$(2,4,6,8)$中二分查找$5$，候选区间如何变化？
        \item 数据经常更新且需范围查询时，散列表是否总是合适？
    \end{itemize}
    \pause
    \begin{block}<2-|handout:0>{参考答案}
        $[0,4)\rightarrow[0,2)\rightarrow[2,2)$，未找到。
        普通散列表不维护顺序；范围查询可考虑有序结构并权衡更新成本。
    \end{block}
\end{frame}
